\chapter{Introdução}

O avanço da tecnologia digital tem transformado profundamente a forma como
indivíduos interagem com sua saúde mental, com seus hábitos cotidianos e com os
recursos disponíveis para o autocuidado. Em um cenário onde a sobrecarga de
informações, a dispersão de atenção e o aumento dos quadros de ansiedade e
depressão se tornam cada vez mais recorrentes, cresce também a busca por
ferramentas que auxiliem no reconhecimento e na gestão das emoções de forma
autônoma e personalizada.

Neste contexto, esta monografia apresenta o desenvolvimento de uma ferramenta
digital voltada ao acompanhamento emocional, com base em princípios de
usabilidade, acessibilidade e respeito à privacidade do usuário. O projeto
parte de duas premissas básicas, uma, de que o nível de consciência afetiva
auxilia no processo de identificação de flutuações de humor, seja períodos de
alta (mania) ou de baixa (depressão), e o segundo elemento é que o registro
pessoal dessas flutuações, uma prática historicamente associada à escrita
íntima e reflexiva, pode ser potencializado com o uso de tecnologias que
favoreçam a auto expressão, o reconhecimento de padrões emocionais e a
construção de uma consciência afetiva mais integrada, sempre que respeite a
privacidade e tenha um cuidado extra com esses tipos de dados sensíveis.

A pesquisa envolve uma fundamentação teórica sobre práticas de auto registro,
saúde mental digital e design de sistemas sensíveis ao contexto do usuário. A
metodologia adotada combina levantamento bibliográfico, aplicação de
questionários e desenvolvimento incremental com base em princípios da
engenharia de software e de desenvolvimento de uma interface centrada no
usuário.
